\documentclass[10pt]{article}

% --- Page Setup ---
% \usepackage[a4paper, margin=0.5cm]{geometry}
\usepackage[a4paper, left=0.5cm, right=0.5cm, top=0.3cm, bottom=0.3cm]{geometry}
\usepackage{multicol}

% \allowdisplaybreaks

\usepackage{booktabs} % Required for professional horizontal rules (\toprule, \midrule, \bottomrule)
\usepackage{array}    % Required for better column formatting

% --- Math and Code ---
\usepackage{amsmath, amssymb}
\usepackage{listings}

% --- Spacing Adjustments ---
\usepackage{enumitem}
\usepackage{titlesec}

\usepackage{nicematrix, tikz}

\usepackage{graphicx}

\usepackage{float}

\usepackage[hidelinks]{hyperref} % hidelinks removes ugly red boxes around links

% 1. Remove space between list items
\setlist{nosep, leftmargin=*} 

% 2. Shrink space around section headers
\titlespacing*{\section}{0pt}{0pt}{0pt}
\titlespacing*{\subsection}{0pt}{1ex plus .2ex minus .2ex}{0.5ex plus .1ex}

% --- NEW: Reduce header font sizes ---
\titleformat*{\section}{\normalsize\bfseries\color{blue!75!black}}
\titleformat*{\subsection}{\small\bfseries\color{magenta}}

% 3. Add a thin line between columns for readability
\setlength{\columnseprule}{0.4pt}
\setlength{\columnsep}{1.5em}

% --- Custom Code Formatting ---
\lstset{
    basicstyle=\ttfamily\scriptsize,
    breaklines=true,
    aboveskip=2pt,
    belowskip=2pt
}

\usepackage{xcolor}
% \newcommand{\sectionbox}[1]{%
%      \colorbox{blue!75!black}{\color{white}\small\bfseries~#1~}%
%    }
%    \titleformat{\section}{}{}{0em}{\sectionbox}
%    \titleformat*{\subsection}{\normalsize\bfseries}

\raggedbottom

\setlength{\parskip}{0pt}
\setlength{\baselineskip}{0pt plus 0.2pt} 


\begin{document}

\section*{Interfaces}


Analog signal is continuous (real world), digital is discrete (for processing).\\
\begin{itemize}
    \item Analog signal is equivalent to a digital signal with infinitely small sampling interval.\\
    \item Nyquist Theorem: For digital signal to be able to represent its analog equiv adequately, minimum sampling rate is 2x the freq of analog signal
\end{itemize}
% Analog-to-digital (ADC) \& Digital-to-analog (DAC) convert the signals for processing.
Sequence of processing real-world data:
\begin{itemize}
    \item ADC $\to$ Digital Processor $\to$ DAC $\to$ smoothing filter.
\end{itemize}
\begin{center}
    \includegraphics[width=0.3\textwidth]{images/adc-dac.png}
\end{center}
Interface is a boundary where 2 or more devices meet to exchange info.\\
Considerations when interfacing devices:
\begin{itemize}
    \item Electrical signal level: output voltage level should not exceed the max allowable input voltage of the input.
    \item Check compatibility: $V_{OH} > V_{IH}$ and $V_{OL} < V_{IL}$.
\end{itemize}
Signals transmitted by differential signals have better noise tolerance.
\begin{center}
    \includegraphics[width=0.35\textwidth]{images/single-vs-differential.png}
\end{center}
\begin{center}
    \includegraphics[width=0.4\textwidth]{images/differential-signal.png}
\end{center}
% \textbf{\textcolor{magenta}{Parallel data trf}}\\
\subsection*{Parallel data trf}


SRAM, DRAM, ROM uses parallel data trf. 
Synchronous in nature with some strobe signal to inform receiver when to latch in the data.\\
Problems:
\begin{itemize}
    \item Signal skew: Signals arrive at diff times 
    \begin{itemize}
        \item Because of variation in propagation delay between signals from the same data bus.
        \item Caused by difference in resistance and capacitance of different wires, so time taken for voltage change on wire to reach steady state differs (time constant $\tau=RC$).\\
        \item Could be further caused by variation in PCB trace length/width (serpentine/trombone).
    \end{itemize}
    \begin{center}
        \includegraphics[width=0.4\textwidth]{images/signal-skew.png}
    \end{center}
    \item Crosstalk: Undesired coupling of signals from one circuit to another
    \begin{itemize}
        \item  Beacuse of close placement of data lines.
        \item  Interferance could be via electromagnetic domain (act like antenna) or electrical domain.
    \end{itemize}
    \item Both signal skew and crosstalk limit: the max no. of lines on the bus, max clock rate, max no. of bytes trfed in 1 block, min delay bw each byte trf.
\end{itemize}
So parallel data trf has faster data trf rate, but susceptible to signal skew and crosstalk, and higher hardware cost and bulkier.\\
\subsection*{Serial data trf}


Data trf one bit a time, so trf rate lower compared to parallel trf, and hardware interface design more complex as it needs to handle serial to parallel conversion (processor process in bytes).\\
But able to trf data reliably over longer distance with lower cost, and less affected by signal skew and crosstalk.\\
\textbf{Serial data trf modes}
\begin{itemize}
    \item Simplex: data trf in one direction
    \item Half-duplex: data trf in both directn, but RX/TX mutually exclusive
    \item Full duplex: 2 data lines, simultaneous RX/TX
\end{itemize}
\begin{center}
    \includegraphics[width=0.35\textwidth]{images/simplex-duplex.png}
\end{center}
\textbf{Nature of communication}
\begin{itemize}
    \item Synchronous: common clock signal bw TX and RX
    
    
    \begin{itemize}
        \item I2C: not tested
        \item SATA: not tested
        \item Serial Peripheral Interface (SPI) Bus.
        \begin{itemize}
            \item Initiate trf: master pulls the Slave Select (SS) low to enable the target SPI slave.
            \item Data on the Master-Out-Slave-In (MOSI) and MISO are latched on the rising edge of the SCLK signal.
            \item Multiple slaves can be connected to one master.
        \end{itemize}
        \begin{center}
            \includegraphics[width=0.4\textwidth]{images/spi-bus.png}
        \end{center}
    \end{itemize}
    \item Asynchronous: no common clock, have to agree on pre-fixede clock freq
    
    
    \begin{itemize}
        \item Universal Asynchronous Receiver Transmitter (UART)
        \begin{itemize}
            \item IDLE = logic 1.
            \item TX sends START bit (logic 0). Receiver detects the start of transmission.
            \item Then DATA is sent at some Baud rate  = no. of bits/second.
            \item PARITY bit may be sent for error checking.
            \item STOP bit (logic 1) terminates the transmission.
            \item The transmission configuration will follow DATA:PARITY:STOP. E.g. 7O1 means 7 data bits, Odd parity, 1 Stop bit.\\
            Parity can be Odd, Even, or No parity. Stop bit can be 1 or 2.
            \item LSB is sent first.
            \item Time for each bit $T = \frac{1}{\text{Baud Rate}}$.
        \end{itemize}
        \begin{center}
            \includegraphics[width=0.4\textwidth]{images/uart.png}
        \end{center}
        \item Internal clock of UART RX will typically run a multiple times (4X here) of the baud rate, to time the sampling close to the middle of each data bit.
        \begin{center}
            \includegraphics[width=0.3\textwidth]{images/uart-rx.png}
        \end{center}
        \item If RX samples at different baud rate, the result will be different.
        \begin{center}
            \includegraphics[width=0.4\textwidth]{images/different-baud.png}
        \end{center}
    \end{itemize}
\end{itemize}
\subsection*{Wired}


\textbf{USB}


\begin{itemize}
    \item Tiered-star topology: USB host (usually the computer) connected to peripherals (e.g. mouse) directly or through USB hub.
    \begin{center}
        \includegraphics[width=0.4\textwidth]{images/usb-topology.png}
    \end{center}
    \item Upon connection, the host and device have to undergo enumeration (exchange info on capability \& requirement), then the host checks for device driver to support the USB device.
\end{itemize}
\textbf{High-definition Multimedia Interface (HDMI)}


\begin{itemize}
    \item Audio/video interface for transmitting uncompressed video/audio data.
\end{itemize}


\subsection*{Wireless}
Industrial, Scientific, Medical (ISM) radio freq (RF) bands are free-to-use frequency bands. The 2.4Ghz and 5.8Ghz bands are standard world wide.\\
\textbf{WiFi}


\begin{itemize}
    \item Operates in 2.4Ghz and 5.8Ghz RF range.
    \item Infrastructure mode: uses Star topology, at the center of the network is an Access Point or Router, connected to devices at the end.
    \item Adhoc mode: Peer to Peer connection.
    \item Transmission range generally 20m-150m.
\end{itemize}
\textbf{Bluetooth}


\begin{itemize}
    \item Mainly for low data rate wireless transmission, focus on low power consumption.
    \item Operates in 2.4Ghz range.
    \item Star topology.
    \item Tranmission range up to 100m but kept to 10-20m for lower power consump.
    \item 2 protocols: BT classic and BLE (Low energy), but most BT hosts can connect to both.
\end{itemize}
\textbf{Factors affecting transmission range}


\begin{itemize}
    \item Transmissn power: higher means greater range.
    \item Transmissn freq: higher freq, higher attenuation, lower range
    \item Interferance between closer transmitters in the same freq band.
    \begin{itemize}
        \item To mitigate, the freq band is divided into channels, and each device uses different channels.
        \item Another solution is frequency hopping from one channel to another, which is what BT uses.
    \end{itemize}
\end{itemize}

\end{document}
